\chapter{Introduzione}
\section{Presentazione del progetto}
Sapienza-DC~\cite{sapienzaDC_repo} è stato realizzato come progetto finale per il corso di Laboratorio di Ingegneria Informatica.

In seguito alla realizzazione di un progetto preliminare, il quale consisteva nella realizzazione di un parser per opportuni domini assegnati a ciascun gruppo partecipante, sono stati selezionati i migliori progetti, al fine di avere l'opportunità di accedere ad un progetto personalizzato.

Tra i vari progetti personalizzati veniva proposta la realizzazione di un chatbot di dominio capace di ricevere in input una query da parte dell'utente, selezionare un dominio su cui cercare informazioni per rispondere alla domanda, ricercare fonti su pagine web appartenenti a quest'ultimo, estrarre e processare i contenuti rilevanti da fornire a un LLM\footnote{LLM - \textit{Large Language Model}: un modello linguistico addestrato su ingenti quantità di dati per la comprensione, l'analisi e la generazione di linguaggio naturale.} avente il compito di rispondere al quesito posto, fornendo aggiuntivamente le fonti di riferimento impiegate.

\section{Requisiti di progetto}
Le specifiche fornite preliminarmente per la realizzazione del progetto consistevano in una descrizione generale dell'architettura della pipeline da implementare:

\begin{enumerate}[label=(\roman*)]
    \item consentire all'utente di porre al proprio sistema una query;
    \item assicurarsi che la query del'utente possa essere risolta sfruttando i domini disponibili;
    \item ricercare online pagine (nei domini disponibili) che possano contenere informazioni utili;
    \item parsare le pagine (sfruttando i parser dell'esonero) e fornire ad un LLM la query assieme ai dati estratti da queste ultime;
    \item restituire la risposta all'utente.
\end{enumerate}

\noindent
In aggiunta a quanto appena elencato, si richievano inoltre diverse funzionalità aggiuntive. In particolare, il sistema implementato sarebbe dovuto essere capace di: gestire richieste ambigue, chiedendo, ove necessario, chiarimenti all'utente; rigettare la query dell'utente qualora quest'ultima fosse ritenuta inaccettabile o inopportuna; classificare i contenuti delle pagine web per grado di pertinenza alla domanda, valutando eventualmente la possibilità di estrarre dati da domini qualunque in caso di query esterna ai domini assegnati; riassumere o selezionare porzioni delle informazioni ottenute in fase di parsing per evitare di saturare la \textit{context window}\footnote{context window: anche detta ``finestra di contesto'', è la massima quantità di testo, misurata in token, che un LLM può analizzare in una singola sessione.} del large language model; citare le fonti e stimare l'affidabilità della risposta fornita, gestendo opportunamente il caso in cui le informazioni in possesso fossero parziali o insufficienti per rispondere alla domanda dell'utente.

Di seguito viene proposto il diagramma utilizzato per presentare i requisiti di progetto.

\begin{figure}[htbp]
    \centering
    \includegraphics[width=0.9\textwidth]{diagramma_requisiti_progetto}
    \caption{Diagramma esemplificativo fornito assieme ai requisiti di progetto.}
    \label{fig:requisiti_progetto}
\end{figure}

\noindent
Si noti come tali specifiche rappresentino solo una linea guida alla realizzazione del progetto, descrivendo solo gli aspetti generali e le funzionalità fondamentali che quest'ultimo dovrebbe implementare.

Questo è voluto: tra gli obiettivi del progetto personalizzato vi è proprio quello di lasciare al gruppo la libertà, e dunque anche le difficoltà che ne conseguono, di produrre quanto richiesto utilizzando gli strumenti nonché tecnologie ritenute più consone.

A tal riguardo, nel capitolo che segue introdurremo le principali tecnologie e framework utilizzati per sviluppare Sapienza-DC.   

